\section*{I. Tóm tắt Thực thi}

Báo cáo này cung cấp phân tích định lượng về quy trình tối ưu hóa mạng tích hợp Không gian - Trên không - Mặt đất hỗ trợ Đa truy nhập phi trực giao (SAGIN-NOMA). Mục tiêu cốt lõi của nghiên cứu là cực tiểu hóa xác suất dừng toàn phần ($P_{\text{out}}$) của hệ thống, nâng cao độ tin cậy dịch vụ trên kênh fading Rician tổng quát, đồng thời bảo đảm tính khả thi nghiêm ngặt của kỹ thuật Khử nhiễu liên tiếp (SIC). Báo cáo trình bày chi tiết mô hình hệ thống hai chặng (Vệ tinh $\to$ Cụm tiếp sóng UAV với cơ chế PRS $\to$ Hai người dùng mặt đất), dẫn giải giải tích dạng đóng chính xác qua hàm Marcum-Q, cùng hai phương pháp tối ưu hóa: (i) Tối ưu hóa Luân phiên (Alternating Optimization - AO) số học đóng vai trò nghiệm tham chiếu chuẩn (Benchmark); và (ii) Tác tử Học tăng cường sâu (Deep Reinforcement Learning - DRL) tùy biến theo kiến trúc Actor-Critic cho phép ra quyết định điều khiển quỹ đạo bay và phân bổ công suất trực tuyến ở cấp độ mili-giây. Các kết quả thực nghiệm mô phỏng, phân tích so sánh định lượng và các thông số kỹ thuật triển khai thực tế trên các nền tảng bay 6G được trình bày chi tiết trong báo cáo này.
